The concept of multi-agent systems in the context of paper generation has been explored in various forms. \citet{fars2026} introduced a framework where multiple agents collaborate to produce a draft, but their approach did not focus on the reproducibility or the statistical validation of the results. \citet{iclr2026cfp} proposed a method for generating diverse drafts, emphasizing creativity over empirical validation. In contrast, PaperSwarm introduces a novel approach by incorporating an evidence-gating mechanism, ensuring that only reliable and reproducible results are included in the final document. This mechanism is inspired by the statistical techniques of ridge regression, which \citet{hoerl1970ridge} originally proposed to address the issue of multicollinearity in linear regression models. The mean test MSE of the minimum-norm OLS solution, averaged over \result{cl-seeds} seeds, was \result{cl-ols-mean}, with a high standard deviation of \result{cl-ols-std} across seeds. In comparison, the mean test MSE of ridge regression at the selected penalty \result{cl-alpha}, which was \result{cl-ridge-mean}, showed lower variability, with a standard deviation of \result{cl-ridge-std}. This relative reduction in mean test MSE achieved by ridge over OLS was \result{cl-improve} percent, indicating the robustness and reliability of the selected penalty. The irreducible noise floor, used as an MSE reference, was set to \result{cl-noise}, and the design matrix contained \result{cl-features} features with \result{cl-train} training samples per split and \result{cl-test} test samples per split. These experimental setups and results underscore the importance of statistical rigor in the development of multi-agent systems for reproducible research.
